
Or \(\lambda^p(\xi,p^*(i_*(y)))-\lambda^p(1,p^*(i_*(y)))\) est divisible par \(1-\xi\) :
pour le voir, on peut se placer dans le \(\lambda\)-anneau universel qui sert à définir les \(\lambda^p(\xi,X)\) (V 5); alors \(\lambda^p(\xi,X)-\lambda^p(1,X)\) est un polynôme en \(\xi\), nul pour \(\xi=1\), donc divisible par \(\xi-1\). Par ailleurs, l'équation de dépendance de \(\xi\) sur \(K^{\bullet}(X)\) est \((1-\xi)^2=0\) (VI (1.13.1)). On obtient donc :
\[
\begin{aligned}
(s\circ i)_*\bigl(\lambda^p(N+1,y)\bigr)
&\equiv \lambda^p(1,p^*(i_*(y)))(1-\xi) \quad\text{mod.\ torsion}\\
&\equiv p^*(\lambda^p(1,i_*(y)))(1-\xi) \quad\text{mod.\ torsion}.
\end{aligned}
\]
En appliquant \(p_*\), on obtient :
\begin{equation}\tag{4.3.4}
  i_*\bigl(\lambda^p(N+1,y)\bigr)
  \equiv \lambda^p(1,i_*(y)) \quad\text{mod.\ torsion},
\end{equation}
car \(p_*(\xi)=0\).

D'après la formule de (V 5.6), appliquée à \(L=1\), \(\lambda^p(N+1,y)\) est le coefficient de \(t^p\) dans la série :
\[
  \frac{\sum_{k\ge1}\lambda^k(N,y)kt^k}
       {1+\sum_{k\ge1}\lambda^k(N,y)\lambda_{-1}(N)t^k}.
\]

De même, la même formule appliquée à \(N=0\) et \(L=1\) montre que \(\lambda^p(1,y)\) est le coefficient de \(t^p\) dans la série :
\[
  \frac{\sum_{k\ge1}\lambda^k(y)kt^k}
       {\sum_{k\ge0}\lambda^k(y)t^k}.
\]

Les congruences (4.3.4) peuvent donc s'écrire :
\[
i_*\left(
  \frac{\sum_{k\ge1}k\lambda^k(N,y)t^k}
       {1+\sum_{k\ge1}\lambda^k(N,y)\lambda_{-1}(N)t^k}
\right)
\equiv
  \frac{\sum_{k\ge1}k\lambda^k(i_*(y))t^k}
       {\sum_{k\ge0}\lambda^k(i_*(y))t^k}
       \quad\text{mod.\ torsion}
\]
